\documentclass[12pt,a4paper]{report}
\usepackage[utf8]{inputenc}
\usepackage[margin=1in]{geometry}
\usepackage{fancyhdr}
\usepackage{graphicx}
\usepackage{listings}
\usepackage{xcolor}
\usepackage{amsmath}
\usepackage{amssymb}
\usepackage{array}
\usepackage{tabularx}
\usepackage{hyperref}
\usepackage{float}

% Code styling
\lstset{
    language=Python,
    basicstyle=\ttfamily\small,
    keywordstyle=\color{blue},
    commentstyle=\color{gray},
    stringstyle=\color{red},
    breaklines=true,
    showstringspaces=false,
    tabsize=4,
    frame=single,
    rulecolor=\color{black},
    backgroundcolor=\color{white}
}

% Header and Footer
\pagestyle{fancy}
\fancyhf{}
\lhead{NEXUS Compiler}
\rhead{Compiler Design CS204}
\cfoot{\thepage}

\title{\textbf{NEXUS Language Compiler}\\
\vspace{0.5cm}
\textit{Grammar Specification and Language Semantics}\\
\vspace{1cm}
\large Compiler Design (CS204)}

\author{
\textbf{Indian Institute of Information Technology (IIIT)}\\
\textbf{Kurnool, Andhra Pradesh}\\
\vspace{1cm}
\textbf{Submitted By:}\\
Anjani Kumar Singh \quad (524CS0003)\\
Poojith \quad (124CS0029)\\
\vspace{1cm}
\textbf{Supervised By:}\\
Prof. Nagaraju
}

\date{\today}

\begin{document}

\maketitle

\tableofcontents
\newpage

\chapter{Introduction}

\section{Overview}
The NEXUS Language is a modern, imperative programming language designed with simplicity and expressiveness in mind. This compiler project implements a complete compilation pipeline that transforms NEXUS source code into multiple target languages including Python, C, and x86-64 Assembly.

\section{Compilation Phases}
The NEXUS compiler implements the following phases:

\begin{enumerate}
    \item \textbf{Lexical Analysis:} Tokenization of source code
    \item \textbf{Syntax Analysis:} Abstract Syntax Tree (AST) generation
    \item \textbf{Semantic Analysis:} Type checking and symbol validation
    \item \textbf{Intermediate Code Generation:} Three-Address Code (TAC) generation
    \item \textbf{Code Optimization:} TAC optimization with multiple strategies
    \item \textbf{Code Generation:} Target code generation (Python, C, Assembly)
\end{enumerate}

\chapter{Language Specification}

\section{Lexical Elements}

\subsection{Keywords}
The NEXUS language defines the following reserved keywords:

\begin{center}
\begin{tabular}{|c|c|c|}
\hline
\textbf{Keyword} & \textbf{Purpose} & \textbf{Example} \\
\hline
\texttt{hold} & Variable declaration & \texttt{hold x = 5;} \\
\hline
\texttt{show} & Print statement & \texttt{show x;} \\
\hline
\texttt{when} & Conditional statement (if) & \texttt{when (x > 5) \{ ... \}} \\
\hline
\texttt{otherwise} & Else branch & \texttt{otherwise \{ ... \}} \\
\hline
\end{tabular}
\end{center}

\subsection{Data Types}
The NEXUS language supports the following primitive data types:

\begin{itemize}
    \item \textbf{int:} Integer values (default type)
    \item \textbf{string:} Text literals
    \item \textbf{unknown:} Type inference for expressions
\end{itemize}

\subsection{Operators}

\subsubsection{Arithmetic Operators}
\begin{center}
\begin{tabular}{|c|c|c|}
\hline
\textbf{Operator} & \textbf{Symbol} & \textbf{Precedence} \\
\hline
Addition & \texttt{+} & 2 \\
\hline
Subtraction & \texttt{-} & 2 \\
\hline
Multiplication & \texttt{*} & 1 \\
\hline
\end{tabular}
\end{center}

\subsubsection{Relational Operators}
\begin{center}
\begin{tabular}{|c|c|}
\hline
\textbf{Operator} & \textbf{Symbol} \\
\hline
Greater Than & \texttt{>} \\
\hline
Less Than & \texttt{<} \\
\hline
Greater Than or Equal & \texttt{>=} \\
\hline
Less Than or Equal & \texttt{<=} \\
\hline
Equal To & \texttt{==} \\
\hline
Not Equal To & \texttt{!=} \\
\hline
\end{tabular}
\end{center}

\subsubsection{Assignment Operator}
\begin{itemize}
    \item \textbf{Assignment:} \texttt{=}
\end{itemize}

\section{Terminals and Non-Terminals}

\subsection{Terminals (Tokens)}
Terminals are the basic symbols that form the vocabulary of the language:

\begin{center}
\begin{tabular}{|l|l|l|}
\hline
\textbf{Token Type} & \textbf{Pattern} & \textbf{Example} \\
\hline
IDENTIFIER & \texttt{[a-zA-Z\_][a-zA-Z0-9\_]*} & \texttt{variable, expr, x} \\
\hline
INTEGER & \texttt{[0-9]+} & \texttt{5, 100, 2025} \\
\hline
STRING & \texttt{"[^"]*"} & \texttt{"Hello World"} \\
\hline
VAR & \texttt{hold} & \texttt{hold} \\
\hline
PRINT & \texttt{show} & \texttt{show} \\
\hline
IF & \texttt{when} & \texttt{when} \\
\hline
ELSE & \texttt{otherwise} & \texttt{otherwise} \\
\hline
ASSIGN & \texttt{=} & \texttt{=} \\
\hline
PLUS & \texttt{+} & \texttt{+} \\
\hline
MINUS & \texttt{-} & \texttt{-} \\
\hline
MULTIPLY & \texttt{*} & \texttt{*} \\
\hline
LEFT\_PAREN & \texttt{(} & \texttt{(} \\
\hline
RIGHT\_PAREN & \texttt{)} & \texttt{)} \\
\hline
LEFT\_BRACE & \texttt{\{} & \texttt{\{} \\
\hline
RIGHT\_BRACE & \texttt{\}} & \texttt{\}} \\
\hline
SEMICOLON & \texttt{;} & \texttt{;} \\
\hline
COMPARISON & \texttt{>, <, >=, <=, ==, !=} & \texttt{>} \\
\hline
NEWLINE & \texttt{$\backslash$n} & Line break \\
\hline
EOF & End of input & - \\
\hline
\end{tabular}
\end{center}

\chapter{Context-Free Grammar}

\section{BNF Notation}
The following Context-Free Grammar (CFG) defines the NEXUS language syntax in Backus-Naur Form (BNF):

\lstinputlisting[language=bash, caption={NEXUS Grammar in BNF}, label=grammar_bnf]{
grammar_rules.txt
}

\section{Grammar Rules}

\subsection{Program Structure}
\begin{verbatim}
Program         → Statement*
Statement       → VariableDeclaration | IfStatement | PrintStatement
EOF             → (implicit end of input)
\end{verbatim}

\subsection{Variable Declaration}
\begin{verbatim}
VariableDeclaration → VAR IDENTIFIER ASSIGN Expression SEMICOLON
                    | VAR IDENTIFIER ASSIGN Expression NEWLINE
\end{verbatim}

\textbf{Semantics:}
\begin{itemize}
    \item Declares a variable with the given identifier
    \item Initializes it with an expression result
    \item The variable type is inferred from the expression
    \item Duplicate declarations result in semantic errors
\end{itemize}

\subsection{Expressions}
\begin{verbatim}
Expression      → AddSubExpr

AddSubExpr      → MulDivExpr ((PLUS | MINUS) MulDivExpr)*

MulDivExpr      → PrimaryExpr (MULTIPLY PrimaryExpr)*

PrimaryExpr     → LITERAL
                | IDENTIFIER
                | LEFT_PAREN Expression RIGHT_PAREN
\end{verbatim}

\textbf{Operator Precedence (highest to lowest):}
\begin{enumerate}
    \item Parentheses \texttt{( )}
    \item Multiplication \texttt{*}
    \item Addition \texttt{+}, Subtraction \texttt{-}
\end{enumerate}

\subsection{Control Flow Statements}

\subsubsection{If-Else Statement}
\begin{verbatim}
IfStatement     → IF LEFT_PAREN BoolExpr RIGHT_PAREN Block
                  (ELSE Block)?

BoolExpr        → Expression COMPARISON Expression
                | Expression

Block           → LEFT_BRACE Statement* RIGHT_BRACE
\end{verbatim}

\textbf{Semantics:}
\begin{itemize}
    \item Evaluates the condition expression
    \item Executes ThenBranch if condition is true
    \item Executes ElseBranch if condition is false (if provided)
    \item Supports nested if-else statements
    \item Creates new scope for block variables
\end{itemize}

\subsubsection{Print Statement}
\begin{verbatim}
PrintStatement  → PRINT (Expression | STRING) SEMICOLON
\end{verbatim}

\textbf{Semantics:}
\begin{itemize}
    \item Outputs the value of an expression or string literal
    \item Expressions are evaluated before printing
    \item String literals are printed as-is
\end{itemize}

\chapter{Semantic Analysis}

\section{Type System}

\subsection{Type Rules}
\begin{itemize}
    \item \textbf{Integer Literals:} Type is \texttt{int}
    \item \textbf{String Literals:} Type is \texttt{string}
    \item \textbf{Variables:} Type is inferred from initialization
    \item \textbf{Binary Operations:}
    \begin{itemize}
        \item \texttt{int + int} → \texttt{int}
        \item \texttt{int - int} → \texttt{int}
        \item \texttt{int * int} → \texttt{int}
        \item Mixing types results in type error
    \end{itemize}
    \item \textbf{Comparison Operations:} Returns boolean (int: 0 or 1)
\end{itemize}

\section{Symbol Table}

The symbol table maintains information about declared identifiers:

\begin{center}
\begin{tabular}{|l|l|l|}
\hline
\textbf{Attribute} & \textbf{Description} & \textbf{Value} \\
\hline
Name & Identifier name & string \\
\hline
Type & Data type & int, string, unknown \\
\hline
Kind & Variable/Function & variable \\
\hline
Scope & Declaration scope & global, local \\
\hline
\end{tabular}
\end{center}

\section{Scope Management}

\subsection{Scope Levels}
\begin{itemize}
    \item \textbf{Global Scope (Level 0):} Top-level declarations
    \item \textbf{Block Scope (Level > 0):} Variables in if/else blocks
    \item \textbf{Scope Rules:}
    \begin{itemize}
        \item Variables are visible in their scope and nested scopes
        \item Shadowing is allowed (inner scope can hide outer scope)
        \item Duplicate declarations in same scope are not allowed
    \end{itemize}
\end{itemize}

\section{Semantic Checks}

\subsection{Type Checking}
\begin{itemize}
    \item All operands in binary operations must be compatible types
    \item Comparison operators work on int types
    \item Assignment must match variable type
\end{itemize}

\subsection{Name Resolution}
\begin{itemize}
    \item References must resolve to declared identifiers
    \item Undefined variable access is reported as error
    \item Forward references are not allowed
\end{itemize}

\subsection{Control Flow Analysis}
\begin{itemize}
    \item Conditions must be evaluable expressions
    \item Nested control structures are supported
\end{itemize}

\chapter{Intermediate Code Generation}

\section{Three-Address Code (TAC)}

\subsection{TAC Statement Types}
\begin{center}
\begin{tabular}{|l|l|l|}
\hline
\textbf{Type} & \textbf{Format} & \textbf{Example} \\
\hline
Assignment & \texttt{t = a op b} & \texttt{t0 = a + b} \\
\hline
Unary & \texttt{t = x} & \texttt{t1 = a} \\
\hline
Jump & \texttt{JUMP L} & \texttt{JUMP L1} \\
\hline
Conditional Jump & \texttt{JUMP\_IF\_FALSE L} & \texttt{JUMP\_IF\_FALSE L0} \\
\hline
Label & \texttt{L:} & \texttt{L0:} \\
\hline
Output & \texttt{WRITE a} & \texttt{WRITE result} \\
\hline
\end{tabular}
\end{center}

\subsection{Temporary Variables}
\begin{itemize}
    \item Temporary variables are named \texttt{t0}, \texttt{t1}, \texttt{t2}, etc.
    \item Each complex expression generates intermediate temporaries
    \item Register allocation can reuse temporaries after their last use
\end{itemize}

\section{Code Optimization Techniques}

The compiler implements the following optimizations:

\subsection{Dead Code Elimination}
Removes unreachable code and unused variable assignments.

\subsection{Constant Folding}
Evaluates constant expressions at compile time.

Example: \texttt{x = 5 + 3;} → \texttt{x = 8;}

\subsection{Common Subexpression Elimination}
Identifies and reuses results of identical expressions.

Example:
\begin{verbatim}
t0 = a + b
t1 = a + b  (eliminated, reuse t0)
\end{verbatim}

\subsection{Loop Unrolling}
Unrolls simple loops for performance improvement.

\subsection{Peephole Optimization}
Performs local optimizations on instruction sequences.

\subsection{Strength Reduction}
Replaces expensive operations with cheaper ones.

Example: \texttt{x * 2} → \texttt{x + x}

\chapter{Code Generation}

\section{Python Code Generation}

Python code is generated with the following features:
\begin{itemize}
    \item Standard Python 3 syntax
    \item Function wrapper: \texttt{def main()}
    \item Return statement for compatibility
    \item Print statements for output
\end{itemize}

\section{C Code Generation}

C code generation includes:
\begin{itemize}
    \item Proper include headers (\texttt{stdio.h}, \texttt{stdlib.h})
    \item Variable declarations at function start
    \item Standard \texttt{main} function signature
    \item \texttt{printf} for output
\end{itemize}

\section{x86-64 Assembly Generation}

Assembly targets the x86-64 architecture with:
\begin{itemize}
    \item AMD64 calling conventions
    \item Stack setup (prologue/epilogue)
    \item Register allocation using \texttt{rax}, \texttt{rsi}, \texttt{rdi}, etc.
    \item Memory addressing with stack pointer offsets
\end{itemize}

\chapter{Example Compilation}

\section{Simple Expression Example}

\subsubsection{Source Code}
\begin{lstlisting}
hold x = 5;
hold y = 3;
show x + y;
\end{lstlisting}

\subsubsection{Tokens Generated}
\begin{verbatim}
VAR, IDENTIFIER(x), ASSIGN, INTEGER(5), SEMICOLON, ...
\end{verbatim}

\subsubsection{AST Structure}
\begin{verbatim}
Program
  VariableDeclaration 'x' (5)
  VariableDeclaration 'y' (3)
  PrintStatement (x + y)
\end{verbatim}

\subsubsection{TAC Generated}
\begin{verbatim}
ASSIGN t=x a1=5
ASSIGN t=y a1=3
ADD t=t0 a1=x a2=y
WRITE a1=t0
\end{verbatim}

\subsubsection{Generated Python}
\begin{lstlisting}[language=Python]
def main():
    x = 5
    y = 3
    t0 = x + y
    print(f"Result: {t0}")
\end{lstlisting}

\chapter{Conclusion}

The NEXUS compiler demonstrates a complete implementation of the compilation pipeline from lexical analysis through code generation. The language supports fundamental programming constructs with clear semantics and multiple target code generation capabilities.

\section{Key Achievements}

\begin{itemize}
    \item Full lexical, syntactic, and semantic analysis
    \item Intermediate code generation with optimization
    \item Multi-target code generation (Python, C, Assembly)
    \item Proper error handling and reporting
    \item Comprehensive symbol table management
\end{itemize}

\section{Future Enhancements}

\begin{itemize}
    \item Support for loops (\texttt{while}, \texttt{for})
    \item Function definitions and calls
    \item Array and structure types
    \item Advanced type system features
    \item Garbage collection mechanisms
\end{itemize}

\vspace{2cm}

\begin{center}
\textbf{Acknowledgments}\\
\vspace{0.5cm}
We gratefully acknowledge Prof. Nagaraju for his guidance and support throughout this project.
\end{center}

\end{document}
